% ----------------------------------------------------------
% ELEMENTOS PRÉ-TEXTUAIS
% ----------------------------------------------------------
% \pretextual

% ---
% Capa
% ---
\imprimircapa
% ---

% ---
% Folha de rosto
% (o * indica que haverá a ficha bibliográfica)
% ---
%\imprimirfolhaderosto*
\imprimirfolhaderosto
% ---



% ---
% Agradecimentos
% ---
\begin{agradecimentos}
Agradeço aos meus pais, Simone e Joel, pelo amor incondicional, pela educação que me
formou como pessoa e pelo apoio em cada etapa da minha trajetória acadêmica. Sem a base
que construíram, este trabalho não seria possível.

À minha namorada Karolline, pelo companheirismo constante, pela paciência nos momentos
de ausência e pelo incentivo que nunca faltou nos momentos mais difíceis deste trabalho.
Sua presença tornou essa jornada mais leve e significativa.

Ao meu orientador, pela disponibilidade, pelas orientações precisas e pela confiança
depositada neste trabalho.

Aos professores do curso de Engenharia da Computação da UNISINOS, pelos ensinamentos que
fundamentam esta e todas as demais etapas da minha formação.

A todos que, de alguma forma, contribuíram para que eu chegasse até aqui, meu sincero
agradecimento.
\end{agradecimentos}
% ---


% ---
% RESUMOS
% ---

% resumo em português
\setlength{\absparsep}{18pt} % ajusta o espaçamento dos parágrafos do resumo
\begin{resumo}
Investigou-se a caracterização de padrões de escoamento bifásico com isobutano em
microtubo horizontal de 1\,mm por meio de sinais de intensidade extraídos de panoramas
de imagens. Analisaram-se 26 ensaios de cinco padrões, formando 87 pares do mesmo padrão
e 238 de padrões diferentes. Extraiu-se a média das linhas 18 a 71 dos panoramas,
seguida de remoção da tendência linear e normalização. Aplicou-se a transformada
\textit{wavelet} contínua de Morlet às funções de correlação, com 200 comprimentos de onda
entre 50 e 3.000 pixels, para obter os espectros de potência e cruzado e a coerência
alternativa. Compararam-se a forma dos perfis espectrais dos panoramas completos, pelo
coeficiente de Pearson $r$ dos resíduos logarítmicos, e os máximos de coerência dos
prefixos comuns de cada par, por um escore ponderado. O coeficiente de Pearson $r$
apresentou média de $0{,}8697$ dentro dos padrões e de $0{,}8038$ entre padrões,
indicando maior semelhança espectral média interna, acompanhada de sobreposição das
distribuições. O escore de coerência apresentou comportamento oposto, com médias de
$0{,}1877$ e $0{,}2635$, respectivamente. A auditoria identificou ressalvas de reconstrução
e acúmulo no teto de intensidade em um ensaio anular; sua exclusão exploratória preservou a diferença positiva
entre as médias do coeficiente de Pearson $r$, mas reduziu sua magnitude. Os resultados
permaneceram descritivos, pois os pares compartilharam aquisições e os grupos foram
desbalanceados. A comparação dos perfis constituiu um indício promissor para análises
adicionais, sem demonstração de significância estatística ou de acurácia de classificação.

 \textbf{Palavras-chave}: escoamento bifásico. microcanais. transformada \textit{wavelet} contínua. coerência \textit{wavelet}. domínio espacial. reconstrução panorâmica. identificação de padrões de escoamento.
\end{resumo}

% resumo em inglês
\begin{resumo}[Abstract]
 \begin{otherlanguage*}{english}
Two-phase flow patterns of isobutane in a horizontal 1\,mm microtube were investigated
using intensity signals extracted from image panoramas. Twenty-six acquisitions from
five patterns were analyzed, yielding 87 within-pattern and 238 between-pattern pairs.
The mean intensity of rows 18 through 71 was extracted, followed by linear detrending
and normalization. The Morlet continuous wavelet transform was applied to correlation
functions at 200 wavelengths from 50 to 3,000 pixels to obtain the power and cross
spectra and alternative coherence. Spectral shape from full panoramas was compared
using Pearson's correlation coefficient $r$ of logarithmic residuals, while coherence
maxima from the common prefixes of each pair were summarized by a weighted score.
Mean Pearson's correlation coefficient $r$ was 0.8697 within patterns and 0.8038 between
patterns, indicating greater mean within-pattern spectral similarity, with overlapping
distributions. The coherence score showed the opposite behavior, with means of 0.1877
and 0.2635, respectively. An audit identified reconstruction concerns and intensity ceiling accumulation
in one annular acquisition; its exploratory exclusion preserved the positive difference
between mean Pearson's correlation coefficients $r$, while reducing its magnitude.
The findings remained descriptive because pairs shared acquisitions and groups were
imbalanced. Spectral profile comparison provided a promising indication for further
analysis, without demonstrating statistical significance or classification accuracy.

   \vspace{\onelineskip}

   \noindent
   \textbf{Keywords}: two-phase flow. microchannels. continuous wavelet transform. wavelet coherence. spatial domain. panorama reconstruction. flow pattern identification.
 \end{otherlanguage*}
\end{resumo}


% ---
% inserir lista de ilustrações
% ---
\pdfbookmark[0]{\listfigurename}{lof}
\listoffigures*
\cleardoublepage
% ---

% ---
% inserir lista de quadros
% ---
\pdfbookmark[0]{\listofquadrosname}{loq}
\listofquadros*
\cleardoublepage
% ---

% ---
% inserir lista de tabelas
% ---
\pdfbookmark[0]{\listtablename}{lot}
\listoftables*
\cleardoublepage
% ---

% ---
% inserir lista de abreviaturas e siglas
% ---
\begin{siglas}
  \item[CNN] Rede Neural Convolucional (\textit{Convolutional Neural Network})
  \item[CWT] Transformada Wavelet Contínua (\textit{Continuous Wavelet Transform})
  \item[DWT] Transformada Wavelet Discreta (\textit{Discrete Wavelet Transform})
  \item[FFT] Transformada Rápida de Fourier (\textit{Fast Fourier Transform})
  \item[LDR] Resistor Dependente de Luz (\textit{Light Dependent Resistor})
  \item[LED] Diodo Emissor de Luz (\textit{Light Emitting Diode})
  \item[LETEF] Laboratório de Estudos Térmicos e Fluidodinâmicos
  \item[LWE] Coeficiente de Energia Wavelet Local (\textit{Local Wavelet Energy coefficient})
  \item[ML] Aprendizado de Máquina (\textit{Machine Learning})
  \item[OS] Sensor Óptico (\textit{Optical Sensor})
  \item[PD] Transdutor de Pressão Diferencial
  \item[PH] Pré-aquecedor (\textit{Pre-Heater})
  \item[PSD] Densidade Espectral de Potência (\textit{Power Spectral Density})
  \item[SNR] Relação Sinal-Ruído (\textit{Signal-to-Noise Ratio})
  \item[SVM] Máquina de Vetores de Suporte (\textit{Support Vector Machine})
  \item[TCC] Trabalho de Conclusão de Curso
  \item[TS] Seção de Teste (\textit{Test Section})
  \item[VS] Seção de Visualização (\textit{Visualization Section})
  \item[WCO] Coerência Wavelet (\textit{Wavelet Coherence})
  \item[WCS] Espectro Cruzado Wavelet (\textit{Wavelet Cross Spectrum})
  \item[WPS] Espectro de Potência Wavelet (\textit{Wavelet Power Spectrum})
\end{siglas}
% ---

% ---
% inserir lista de símbolos
% ---
\begin{simbolos}
  % --- Quantidades espectrais ---
  \item[$W_x(a,b)$] Coeficiente da CWT do sinal $x$ na escala $a$ e posição $b$
  \item[$P_x(a,b)$] Espectro de potência wavelet clássico do sinal $x$
  \item[$C_{xy}(a,b)$] Espectro cruzado wavelet clássico entre os sinais $x$ e $y$
  \item[$\gamma_{xy}(a,b)$] Coerência wavelet clássica entre os sinais $x$ e $y$
  \item[$R_{xx}(\tau)$] Função de autocorrelação do sinal $x$
  \item[$R_{xy}(\tau)$] Função de correlação cruzada entre os sinais $x$ e $y$
  \item[$\widetilde{P}_x(a,b)$] Campo de potência wavelet alternativo do sinal $x$
  \item[$\widetilde{C}_{xy}(a,b)$] Espectro cruzado wavelet alternativo entre $x$ e $y$
  \item[$\widetilde{\gamma}_{xy}(a,b)$] Coerência wavelet alternativa entre $x$ e $y$
  \item[$\lambda_k$] Comprimento de onda espacial na escala $k$ [px]
  \item[$P_{ik}$] Potência alternativa no atraso nulo do ensaio $i$ na escala $k$
  \item[$M_{ik}$] Potência alternativa média nos atrasos do ensaio $i$ na escala $k$
  \item[$z_{ik}$] Resíduo do perfil logarítmico de potência após remoção da tendência
  \item[$r_{ij}$] Coeficiente de Pearson entre os descritores dos ensaios $i$ e $j$
  \item[$\eta$] Piso relativo de potência para o escore de coerência
  \item[$S_{xy}^{(\eta)}$] Escore ponderado dos máximos de coerência por escala
  \item[$\Delta_h$] Contraste descritivo entre médias intra-padrão e inter-padrão da métrica $h$
  % --- Variáveis termofluidodinâmicas ---
  \item[$\beta$] Parâmetro adimensional de identificação do padrão de escoamento
  \item[$G$] Fluxo mássico [kg\,m$^{-2}$\,s$^{-1}$]
  \item[$q''$] Fluxo de calor [kW\,m$^{-2}$]
  \item[$Co$] Número de confinamento
  \item[$\sigma$] Tensão superficial [N\,m$^{-1}$]
  \item[$\rho_l$] Massa específica do líquido [kg\,m$^{-3}$]
  \item[$\rho_v$] Massa específica do vapor [kg\,m$^{-3}$]
  \item[$D_h$] Diâmetro hidráulico [m]
  \item[$i_{lv}$] Calor latente de vaporização [J\,kg$^{-1}$]
\end{simbolos}
% ---

% ---
% inserir o sumario
% ---
\pdfbookmark[0]{\contentsname}{toc}
\tableofcontents*
\cleardoublepage
% ---
